\chapter{绪论}

\section{课题背景}
医学影像记录人体结构和病理变化，是临床诊断、疾病筛查、治疗规划和医学教学中经常使用的资料。磁共振成像（Magnetic Resonance Imaging，MRI）具有较好的软组织对比度，不使用电离辐射，也支持多序列成像。正因为这些特点，MRI 在脑部结构观察、神经系统疾病分析和医学影像教学中具有较高使用价值。

在实际教学和模型实验中，真实 MRI 数据并不总是容易获取和使用。扫描设备成本较高，病例数据又涉及隐私保护和伦理审查，数据标注也常常需要医学专业知识参与。即使能够获得公开数据，不同设备和扫描协议仍会带来灰度分布差异。对人工智能模型来说，数据规模和数据一致性都会影响训练稳定性；对课堂演示和毕业设计展示来说，仅给出单张二维切片，也很难把三维结构关系讲清楚。

生成式人工智能为医学影像教学和实验提供了新的思路。生成对抗网络（Generative Adversarial Network，GAN）通过生成器和判别器的对抗训练学习数据分布，已被用于图像生成、图像增强和跨模态合成等任务\cite{goodfellow2014gan,skandarani2023gans,khader2023ddpm3d}。StyleGAN 系列模型重视图像质量和潜空间控制。StyleGAN2 改进了生成器结构和正则化方式，生成图像质量也进一步提高\cite{karras2020stylegan2}。StyleGAN2-ADA 通过自适应数据增强缓解有限数据训练中的过拟合问题\cite{karras2020ada}。本文选择这一框架，主要是考虑到 MRI 切片灰度结构相对集中，且训练数据来源有限。

当二维切片可以生成之后，新的问题随之出现：这些切片如何被组织成更容易理解的空间表达。三维高斯泼溅（3D Gaussian Splatting，3DGS）使用可优化的三维高斯基元表示场景，并通过可微渲染生成视图\cite{kerbl20233dgs}。原始 3DGS 主要面向多视角照片重建，而 MRI 切片与普通照片的成像方式并不相同。因此，本课题没有简单套用自然场景重建流程，而是借鉴点云初始化和三维高斯表达思路，将连续 MRI 切片组织起来，并通过平台展示三维结果。由此形成了“二维生成 + 三维表达 + 平台展示”的研究路线。

\section{国内外研究现状}
医学图像生成研究主要服务于数据扩充、隐私保护、跨模态合成和辅助教学。早期方法常通过旋转、缩放、裁剪和灰度扰动来扩充数据。这类方法实现简单，但本质上仍是在改变已有图像，难以产生新的结构变化。深度生成模型出现后，GAN、变分自编码器和扩散模型逐渐被用于医学图像合成。

近年医学图像生成研究更加关注数据稀缺、隐私保护和下游任务可用性。Skandarani 等在多模态医学图像合成实验中指出，不同 GAN 架构在 MRI、CT 和眼底图像上的表现并不一致，视觉真实并不必然带来下游任务提升\cite{skandarani2023gans}。Khader 等研究了扩散模型在三维医学图像生成中的应用，说明高质量合成数据需要同时关注图像外观、解剖合理性和切片一致性\cite{khader2023ddpm3d}。Müller-Franzes 等进一步比较了潜在扩散模型与 GAN 在医学图像合成中的表现，提示生成模型评价不能只依赖单一视觉指标\cite{mullerfranzes2023comparison}。Akbar 等围绕脑肿瘤 MR 合成图像开展分割实验，也强调合成数据需要通过具体任务验证其实际价值\cite{akbar2024syntheticmr}。基于这些研究，本文把平台用途限定为教学演示和工程验证，不把输出结果用于临床诊断。

GAN 方法的发展为本课题提供了二维生成基础。Goodfellow 等提出的 GAN 建立了对抗式生成学习的基本框架\cite{goodfellow2014gan}。StyleGAN 引入基于风格的生成结构和中间潜空间，提高了图像合成质量\cite{karras2019stylegan}。StyleGAN2 在此基础上改进结构，并减少原 StyleGAN 中的伪影问题\cite{karras2020stylegan2}。StyleGAN2-ADA 面向有限数据训练场景，使用自适应数据增强提升训练稳定性\cite{karras2020ada}。在更近的脑影像生成研究中，Tudosiu 等强调了三维脑影像生成对形态保持和结构合理性的要求\cite{tudosiu2024brain}。本课题使用公开 MRI 数据切片，病例来源和结构类型仍有限，因此选用 StyleGAN2-ADA 作为二维 MRI 切片生成框架，并在论文中明确其教学演示边界。

三维展示部分更强调工程适配。传统医学三维展示常使用体绘制、表面重建或体素显示，这些方法适合展示真实体数据。本文的输入还包括生成切片，所以必须处理切片连续性、坐标尺度和灰度一致性。3DGS 从点云出发，并用三维高斯表示空间结构\cite{kerbl20233dgs}。2024 年以来，已有研究把 3DGS 引入内窥镜重建、手术场景重建和合成数据生成等医学相关任务，说明显式高斯表示在医学视觉场景中具有较好的工程拓展性\cite{bonilla2024gaussianpancakes,huang2024endo4dgs,guo2024freesurgs,zeng2024surgical3dgs}。因此，本课题把重点放在工程链路上：完成二维切片到点云的转换，并把点云和连续切片用于 3DGS 训练与展示。

\section{项目研究意义}
本课题的价值不只在于复现某一个模型，更在于把原本分散的实验步骤整理成一个可以运行和展示的平台流程。MRI 数据读取、切片预处理、StyleGAN2 训练、图像生成、点云构建、3DGS 训练和结果展示被放在同一条链路中，脚本输出不再只是零散文件，而可以通过页面、日志和图像结果进行查看。

从医学图像智能处理角度看，本课题探索了 StyleGAN2-ADA 在脑部 MRI 二维切片生成中的使用方式，并把生成切片继续组织为连续序列和三维点云。这样处理后，二维生成结果不再停留在单张图片层面，而可以进入后续三维表达流程。对学习者而言，这种设计更容易说明医学体数据从切片到空间结构的转换过程。

从教学应用角度看，平台采用 FastAPI + React 构建前后端分离界面。用户不需要反复输入复杂命令，就可以在网页中完成数据查看、模型生成、影像分析、任务追踪和三维结果展示。这样的形式更适合人工智能专业医学图像处理方向的课程演示、实验复核和毕业答辩展示。

\section{论文主要内容与章节安排}
围绕上述目标，论文工作主要包括四个部分。

第一部分是 MRI 数据预处理。系统读取 IXI 数据集中 NIfTI 格式的 T1 脑部 MRI 数据，提取轴向切片，并完成百分位归一化、尺寸统一和低信息切片过滤。处理后的数据可用于 StyleGAN2-ADA 训练。

第二部分是二维 MRI 切片生成。系统基于预处理切片训练 StyleGAN2-ADA，并在平台中支持权重加载、随机种子设置、截断参数调节和批量切片生成。

第三部分是面向三维表达的序列组织和 3DGS 训练。系统通过潜空间插值得到连续切片，再依据像素坐标和灰度值写出 PLY 点云。结合规则相机参数后，系统完成 3DGS 训练，并保存 checkpoint、训练后点云和渲染图像。

第四部分是 FastAPI + React 交互平台。后端提供生成、分析、任务、资产、体数据和 AI 助手等接口；前端提供仪表盘、生成工作台、医学影像查看器、报告、可信性评估、点云展示和教学中心等页面。用户可以在统一界面中完成主要实验操作。

本文结构安排如下。第 1 章介绍研究背景、研究现状、研究意义和论文结构。第 2 章介绍 MRI、NIfTI、GAN、StyleGAN2、3DGS、FastAPI、React 和 Vite 等相关技术。第 3 章进行系统需求分析。第 4 章给出系统总体设计。第 5 章说明系统关键模块实现。第 6 章进行测试与结果分析。第 7 章说明本文工作完成情况与后续改进方向。
